\chapter{Scope}\label{chap:scope}

\section*{At a glance}
\begin{itemize}
\item \lead{Goal} State which \code{\#erase} inputs the verification covers, and why: those whose
  verification needs nothing beyond lean4lean \code{master}.
\item \lead{Main results} The scope is the predicate \code{InScope}: a closed Lean term, well
  typed through \code{TrS} in the model of an inductive-free program environment that the eraser's
  view agrees with.
\item \lead{Status} \stProved{} every node: \code{ViewAgrees}, \code{InScope}, and their parts
  \code{TrS} and \code{ProgEnv} (Chapter~\ref{chap:model}). The routing lemma
  \code{collectDeps\_not\_outOfFragment} is in Chapter~\ref{chap:core}.
\item \lead{Nodes} \bpnodes{chap:scope}
\item \lead{Reference} MetaRocq \code{erase\_correct}'s hypotheses \code{wf\_ext},
  \code{Sigma} $\sim$\code{\_ext X}, \code{welltyped}; lean4lean \code{TrEnv'}.
\end{itemize}

\section{What the scope is}\label{sec:scope-summary}

\begin{itemize}
\item \lead{In scope} closed Lean terms built from $\lambda$, application, $\Pi$, sorts,
  \code{let}, metadata and constants, whose constants (transitively, through types and values) are
  axioms, definitions (including \code{unsafe} recursive definitions and mutual blocks), theorems
  or opaques, and which are well typed in lean4lean's \code{IsDefEq} through its translation.
\item \lead{In practice} Church-encoded programs over axiomatized base types, with universe
  polymorphism, proofs, type arguments and \code{let}.
\item \lead{Out of scope} \hl{everything that needs an inductive type}: \code{Nat}, \code{Bool},
  structures, classes, \code{match} on data, structural, well-founded and \code{partial}
  recursion; and literals, projections, \code{Quot}, metavariables.
\item \lead{Benchmarks} none is in scope: every program of \code{benchmarks/} takes or returns
  \code{Nat}, \code{Unit}, \code{List} or \code{Bool}.
\item \lead{Routing} \code{\#erase} sends a program whose closure leaves the fragment (an
  inductive, literal, projection, quotient or metavariable) to the unchanged \code{Meta} path, and
  every other program to the pure path (Definition~\ref{def:Erasure-eraseEntry}); the routing
  lemma puts every in-scope input on the pure path
  (Lemma~\ref{lem:EraseProof-collectDeps_not_outOfFragment}).
\end{itemize}

\section{Evidence from lean4lean \texorpdfstring{\code{master}}{master}}\label{sec:scope-evidence}

\subsection{What the model has}\label{sec:scope-has}

\begin{tabular}{p{0.2\linewidth}p{0.4\linewidth}p{0.32\linewidth}}
\toprule
Feature & In lean4lean \code{master} & Consequence \\
\midrule
Terms & \code{VExpr}: \code{bvar}, \code{sort}, \code{const}, \code{app}, \code{lam},
  \code{forallE} (\leanfile{Lean4Lean/Theory/VExpr.lean}) & $\lambda$, application, $\Pi$, sorts, constants \\
Typing & one judgment \code{IsDefEq} with 13 rules, among them \code{beta}, \code{eta},
  \code{proofIrrel}, \code{extra} (\leanfile{Lean4Lean/Theory/Typing/Basic.lean}) & $\beta$, $\eta$, proof
  irrelevance, level equivalence \\
Definitions & a constant and a defining equation (\code{TrEnv'.defn}) & $\delta$ \\
Mutual and unsafe definitions & \code{TrEnv'.mutualDef} & unsafe recursion \\
Theorems, opaques & a constant with a typed value, no equation & in scope; opaques do not unfold
  (DV-10) \\
Axioms & a constant & in scope (DV-21) \\
\code{let}, \code{mdata} & \code{TrExprS.letE} substitutes the value; metadata is transparent &
  in scope; the source semantics keeps \code{let} (DV-1) \\
\bottomrule
\end{tabular}

\subsection{What the model lacks}\label{sec:scope-lacks}

\begin{tabular}{p{0.2\linewidth}p{0.44\linewidth}p{0.28\linewidth}}
\toprule
Missing & Evidence & Consequence \\
\midrule
Inductive types & \code{VInductDecl.WF} and \code{VEnv.addInduct} are sorry definitions (L1,
  L2); the relation \code{AddInduct} has no constructors, so \code{TrEnv'.induct} never fires &
  out: data types, \code{match}, recursors, \code{partial def} \\
Projections & \code{TrProj := sorry}; every lemma about it is a sorry proof & out \\
Literals & the rule \code{TrExprS.lit} needs constants named \code{Nat}, \code{String},
  inductives in a real Lean environment & out; \code{TrS} drops the rule (DV-2) \\
\code{Quot} & \code{TrEnv'.quot} needs a constant \code{Eq}, an inductive in a real environment &
  out; \code{ProgEnv} has no \code{quot} rule (DV-3) \\
Metavariables & no rule of \code{TrExprS}; \code{VLevel.ofLevel} rejects level metavariables &
  out \\
Normalization & no normalization theorem & partial correctness (DV-17) \\
Unique typing, injectivity & theorems resting on L1--L5, and the sorries L4, L6 themselves &
  inherited (Section~\ref{sec:trust-inherited}) \\
\bottomrule
\end{tabular}

\subsection{How the scope follows}\label{sec:scope-follows}

\begin{itemize}
\item \lead{Missing or sorry rules} inductives, projections and metavariables have no usable
  rule.
\item \lead{Rules over inductives} literals and \code{Quot} need constants that a real Lean
  environment declares as inductives, which \code{ProgEnv} cannot contain.
\item \lead{One condition of ours} \code{ProgEnv} requires \code{ConstantInfo.all} to be what the
  elaborator sets: a declaration lists itself, a block member its block. lean4lean's kernel does
  not check \code{all}; the eraser reads it to find recursive constants.
\item \lead{Sub-environments} a real Lean environment always contains \code{Init}'s inductives, so
  the theorem speaks about a program environment \code{P} that the eraser's view agrees with on
  \code{P}'s names (DV-3).
\end{itemize}

\section{The scope as a predicate}\label{sec:scope-predicate}

\begin{definition}[The view agrees with the program]
  \label{def:EraseProof-ViewAgrees}
  \lean{EraseProof.ViewAgrees}
  \leanok
  \uses{def:Erasure-EnvView}
  \srcloc{proof/EraseProof/Core/Scope.lean}{20}
  \lead{In short} \code{ViewAgrees view P}: \hl{\code{view.find?} returns each declaration of
  \code{P} under its name}.
  \lead{Reference} \code{Sigma} $\sim$\code{\_ext X} (\code{abstract\_env\_ext\_rel}), a
  hypothesis of MetaRocq's \code{erase\_correct}.
\end{definition}

\begin{definition}[The scope]
  \label{def:EraseProof-InScope}
  \lean{EraseProof.InScope}
  \leanok
  \uses{def:EraseProof-ViewAgrees, def:EraseProof-ProgEnv, def:EraseProof-TrS}
  \srcloc{proof/EraseProof/Core/Scope.lean}{26}
  \lead{In short} \code{InScope view e}: \hl{some program environment \code{P} with model
  \code{venv}} agrees with the view, and \code{e} translates in \code{venv}:
  \code{ProgEnv P venv}, \code{ViewAgrees view P}, \code{TrS venv Us [] e e'}.
  \lead{Reference} the hypotheses \code{wf\_ext}, \code{Sigma} $\sim$\code{\_ext X},
  \code{welltyped} of MetaRocq's \code{erase\_correct}.
\end{definition}

\section{Examples}\label{sec:scope-examples}

\begin{tabular}{p{0.46\linewidth}p{0.46\linewidth}}
\toprule
Inside & What it exercises \\
\midrule
\code{mul two three}, \code{pow two three}, \code{pred six} over Church numerals & $\lambda$,
  application, $\delta$, erased type arguments \\
\code{pid}, \code{comp}, \code{pflip} at \code{Sort 2}, \code{Sort 1}, \code{Sort 0} & universe
  polymorphism, instances at \code{Prop} \\
\code{guard hP six} with \code{axiom hP : P} & a proof axiom passed by value (an atom, DV-11) \\
\code{let x := two; add x x} & \code{let} of values, types and proofs \\
\code{succ opaqueTwo} & an opaque, stuck in the source \\
\code{uf one}, \code{uf} unsafe recursive & unsafe recursion, erased to a \code{tFix} \\
\bottomrule
\end{tabular}

\begin{tabular}{p{0.46\linewidth}p{0.46\linewidth}}
\toprule
Outside & Gap \\
\midrule
every benchmark, structures, classes, \code{if}, \code{match} on data & inductives \\
structural, well-founded or \code{partial} recursion & inductives \\
\code{s.1}; \code{5}, \code{"x"} & projections; literals \\
\code{Quot.lift f h (Quot.mk r a)} & \code{Quot} needs \code{Eq} \\
\code{(\_ : CNat)}, \code{@pid} with its level open & metavariables \\
\bottomrule
\end{tabular}

\section*{Caveats}
\begin{itemize}
\item \lead{Caveat} The eraser may still fail on an in-scope input (fuel, oracle failure, name
  collision): no output, no claim.
\item \lead{Caveat} The routing test is syntactic and the scope semantic; only the direction
  "in scope implies routed to the verified path" is claimed.
\end{itemize}
